\documentclass[11pt]{article}
\usepackage[margin=1in]{geometry}
\usepackage{amsmath,amssymb,amsthm}
\usepackage{mathrsfs}
\usepackage{hyperref}
\usepackage{enumitem}
\setlength{\emergencystretch}{3em}

\newtheorem{theorem}{Theorem}
\newtheorem{lemma}[theorem]{Lemma}
\newtheorem{proposition}[theorem]{Proposition}
\newtheorem{corollary}[theorem]{Corollary}
\newtheorem{conjecture}[theorem]{Conjecture}
\theoremstyle{definition}
\newtheorem{definition}[theorem]{Definition}
\theoremstyle{remark}
\newtheorem{remark}[theorem]{Remark}

\newcommand{\R}{\mathbb{R}}
\newcommand{\Z}{\mathbb{Z}}
\newcommand{\N}{\mathbb{N}}
\newcommand{\Q}{\mathbb{Q}}
\newcommand{\ee}{\mathrm{e}}
\newcommand{\dist}{\operatorname{dist}}
\newcommand{\dimH}{\dim_{\mathrm{H}}}
\newcommand{\Lam}{\Lambda}

\title{A spectral gap and zero-dimensional level sets for Diophantine\\
approximation along the convergent denominators of \texorpdfstring{$e$}{e}}
\author{Phase-2 report (AI-assisted research; Lean-kernel-checked core)}
\date{Revised November 18, 2026}

\begin{document}
\maketitle

\begin{abstract}
Let $(Q_j)$ be the convergent denominators of $e$ immediately preceding
the partial quotients $2j$, and put $\Lambda(x)=\limsup_j j\|Q_jx\|$,
where $\|t\|$ is the distance from $t$ to the nearest integer.
For each integer $k\ge1$ we identify a gauge $h_k$ by
$h_k(1/Q_n)=(2k+1)^{-n}$ and linear interpolation. On the circle,
the exact level $\{\Lambda=k/4\}$ has infinite, $\sigma$-finite
$\mathcal H^{h_k}$ measure, and the strict sublevel has zero
$\mathcal H^{h_k}$ measure. Every finite sublevel
$\{\Lambda\le c\}$ has Hausdorff dimension zero. The zero level is
$\mathbb Z+\mathbb Z(e+1)/2$ and the positive spectrum has a gap
below $1/4$. The Lean declarations audited for the classification,
gap and dimension bounds use only Lean's three standard axioms.
For ordinary continued fractions, eventual periodicity of the
partial quotients modulo four has the same Hausdorff dimension
as the fixed alphabet $1\bmod 4$, strictly between $13/20$ and $5/7$.
For $0\le\alpha<1$, prescribed admissible limits of
$j^\alpha\|Q_jx\|$ occur on sets of dimension $1-\alpha$.
Almost every input has at least logarithmically many windows with
at most one new convergent.
\end{abstract}

\section{Introduction}

This note studies the metric theory of Diophantine approximation along a
\emph{sparse, super-exponentially growing} subsequence of denominators: the
selected convergent denominators of $\ee$. It is a spin-off of a formalization
project on the (open) irrationality of $e+\pi$; the classification and gap
theorems from that project underlie the metric results presented here.

\subsection{Setup}

The continued fraction of $\ee$ is $[2;1,2,1,1,4,1,1,6,\dots]$. Let $P_j/Q_j$
be the convergent immediately preceding the partial quotient $2j$; the first
few are $3/1,19/7,193/71,2721/1001$. The denominators satisfy
\begin{equation}\label{eq:rec}
Q_{j+1}=(4j+2)Q_j+Q_{j-1}\quad(j\ge2),\qquad Q_1=1,\ Q_2=7,
\end{equation}
so $Q_{j+1}/Q_j\sim 4j$ and $\log Q_j\sim j\log j$ (super-exponential growth).
Write $\|y\|=\dist(y,\Z)$ and define
\[
\Lam(x):=\limsup_{j\to\infty} j\,\|Q_j x\|\in[0,\infty].
\]
Let $L:=\Z+\Z\frac{\ee+1}{2}$, a countable dense subgroup of $\R$.

\subsection{Relation to prior work}\label{sec:prior}

The \emph{phenomenon} (a spectral gap plus zero-dimensional sub-threshold level
sets) is classical: the Lagrange spectrum below $3$ is countable
(dimensions $0$), Hall's ray gives a half-line, and Moreira and Lima--Moreira
establish phase transitions for spectra of Diophantine approximation
\cite{Moreira,LimaMoreira}. For comparison with continued-fraction characterizations see
\cite{KL,BDS}. In particular, put $\beta=(e+1)/2$. The fractions
$R_j/Q_j$, where $R_j=(P_j+Q_j)/2$, are the convergents of
$\beta=[1;1,6,10,14,\ldots]$: the initial pairs are $(1,1)$ and
$(2,1)$ and the remaining partial quotients are $4j+2$ for $j\ge1$.
The Kraaikamp--Liardet criterion in \cite[p.~106, Lem.~4.1]{BDS} then
gives $\|Q_jx\|>Q_j|Q_j\beta-R_j|/4$ infinitely often for $x\notin L$.
The continued-fraction error formula gives
$jQ_j|Q_j\beta-R_j|\to1/4$, hence $\Lambda(x)\ge1/16$ outside $L$.
Together with $\Lambda=0$ on $L$, this recovers the zero-level
classification. The recurrence proof below gives the sharp bound $1/4$.

For the selected sequence here, Corollary~8 of \cite{BLR}
applies directly since $Q_{j+1}/Q_j\to\infty$. It yields full dimension
for $\{x:\|Q_jx\|\to0\}$. Theorem~7 of that paper treats shrinking
intervals with moving centres. We give the elementary construction
for this recurrence below.

Our results concern a specific sequence with four distinguishing
features: (i) a \emph{subsequence} of convergent denominators (every third),
(ii) the \emph{weight} $j$, (iii) the lattice $L=\Z+\Z\frac{\ee+1}{2}$ (note
$(\ee+1)/2$, not $\ee$), and (iv) the \emph{sharp} constant $1/4$. The weight is
\emph{critical} for the covering argument: it makes the per-parent branching
number uniformly bounded, yielding dimension $0$ for \emph{all} finite level
sets (Theorem~\ref{thm:finite}), whereas without the weight the branching grows
linearly and the argument fails. The unweighted zero-set has full dimension, whereas the zero-set
for the weight $j$ equals the countable group $L$.

\section{Classification and spectral gap}

The following two results were established and Lean-kernel-checked in the
Phase-1 project. The Lean foundation proves both error estimates
$\delta_j\to0$ and $|\delta_j|<1/(2jQ_j)$ for $j\ge1$, where
$\delta_j=Q_j\ee-P_j$. The proofs use the selected-convergent
recurrence and its Hermite integral representation.

\begin{theorem}[Classification, {\cite[Thm.~2]{Phase1}}]\label{thm:class}
For $x\in\R$, $\quad j\|Q_jx\|\to0\iff x\in L.$
\end{theorem}

\begin{theorem}[Gap, {\cite[Cor.~3]{Phase1}}]\label{thm:gap}
If $x\notin L$ and $c<1/4$, then it is \emph{not} the case that
$j\|Q_jx\|\le c$ eventually. Equivalently $\Lam(x)\ge1/4$ for $x\notin L$.
\end{theorem}

\section{The spectral gap (machine-checked)}

\begin{theorem}[Spectral gap]\label{thm:spectral}
For every $x\in\R$, either $\Lam(x)=0$ or $\Lam(x)\ge1/4$. Thus
$\Lam(\R)\subseteq\{0\}\cup[1/4,\infty]$, and $\Lam(x)=0\iff x\in L$.
\end{theorem}

\begin{proof}
If $x\in L$, Theorem~\ref{thm:class} gives $j\|Q_jx\|\to0$, hence $\Lam(x)=0$.
If $x\notin L$, Theorem~\ref{thm:gap} gives $\Lam(x)\ge1/4$. The dichotomy and
the characterization of the zero value follow.
\end{proof}

\begin{remark}
In particular $\Lam(\ee)=0$, since $\ee=2\cdot\frac{\ee+1}{2}-1\in L$; $\ee$ is
an exceptional input for this exponent: its value is zero.
Section~\ref{sec:num} reports finite-window maxima for $B\pi$, $\sqrt2$,
the golden ratio, and random reals.
\end{remark}

\section{Sub-threshold level sets have dimension zero (machine-checked)}

For $c\ge0$ write $E_c:=\{x:\Lam(x)\le c\}$.

\begin{theorem}[Sub-threshold level sets]\label{thm:sub}
If $c<1/4$, then $E_c=L$. In particular $E_c$ is countable and
$\dimH E_c=0$.
\end{theorem}

\begin{proof}
If $x\in L$ then $\Lam(x)=0\le c$, so $L\subseteq E_c$. Conversely if
$x\in E_c$ then $\Lam(x)\le c<1/4$, so by Theorem~\ref{thm:spectral}
$\Lam(x)=0$, hence $x\in L$ by Theorem~\ref{thm:class}. Thus $E_c=L$, which is
countable (image of $\Z^2$), and countable sets have Hausdorff dimension $0$.
\end{proof}

\paragraph{Lean status.} Theorems~\ref{thm:spectral} and~\ref{thm:sub}, together
with the countability of $L$ and $\dimH E_c=0$ for $c<1/4$, are kernel-checked in
\texttt{EpiFormal/SpectralGap.lean} (five declarations, no \texttt{sorryAx}).
The declaration audit reports \texttt{propext}, \texttt{Classical.choice}, and
\texttt{Quot.sound} for the classification and spectral-gap results.
In the existing project environment, \texttt{lake build} and
\texttt{lake env lean Audit.lean} both exited successfully. A fresh
dependency fetch in the separate reproduction directory was blocked by
DNS resolution. The gauge and residue theorems use the proofs in this article.

\section{All finite level sets have dimension zero (machine-checked)}

The sub-threshold result uses the sharp classification. For $c\ge1/4$ the level
set $E_c$ is uncountable, but we show it still has dimension $0$. The entire
covering argument --- factorial growth of $Q_j$, uniform branching bound,
surviving-center recursion, periodicity reduction, Hausdorff-content estimate,
and the countable-union assembly --- is \emph{machine-checked}
(\texttt{EpiFormal/DimHCore.lean} + \texttt{EpiFormal/DimHFinite.lean}, base
axioms only, no \texttt{sorryAx}, and \emph{no} continued-fraction axioms).
The elementary proof below mirrors the formalization.

\begin{lemma}[Factorial growth]\label{lem:factorial}
For $j\ge1$, $Q_j\ge j!$. (Lean: \texttt{DimHCore.\allowbreak Qsel\_ge\_factorial}.)
\end{lemma}

\begin{lemma}[Uniform branching bound, machine-checked]\label{lem:branch}
Fix $t\ge0$, $j\ge1$ and a parent center $c$. The number of child centers
$k'\in[0,Q_{j+1}]$ with $|c/Q_j-k'/Q_{j+1}|\le t/(jQ_j)+t/((j+1)Q_{j+1})$ is at
most $B:=\lfloor16t\rfloor+1$, \emph{uniformly in $j$ and $c$}.
(Lean: \texttt{DimHCore.fiber\_card\_le}; proof by injecting $k'\mapsto
\lfloor k'-A\rfloor$ into $[0,\lfloor16t\rfloor]$, using $Q_{j+1}\le(4j+3)Q_j$.)
\end{lemma}

\begin{theorem}[Finite level sets, machine-checked]\label{thm:finite}
For every $c\in\R$, $\dimH E_c=0$. Consequently
$F:=\{x:\Lam(x)<\infty\}=\bigcup_{n\ge1}E_n$ has $\dimH F=0$.
(Lean: \texttt{DimHFinite.dimH\_levelSet\_finite}; \texttt{\#print axioms}
returns \texttt{[propext, Classical.choice, Quot.sound]} only.)
\end{theorem}

\begin{proof}[Proof sketch (elementary covering)]
Fix $c<\infty$ and $t>c$. Since $\Lam(x)\le c<t$ means $j\|Q_jx\|\le t$
\emph{eventually} (a liminf-type condition, \emph{not} a limsup set), we have
\[
E_c\subseteq\bigcup_{J\ge1}\bigcap_{j\ge J}A_j(t),\qquad
A_j(t):=\{x:\|Q_jx\|\le t/j\}.
\]
Each $A_j(t)$ is a union of $Q_j$ intervals of radius $t/(jQ_j)$ centered at
$k/Q_j$. The intersection $C_J:=\bigcap_{j\ge J}A_j(t)$ is a nested Cantor set.
A level-$j$ interval survives to level $j+1$ only if its center $k/Q_j$ is
within $t/(jQ_j)+t/((j+1)Q_{j+1})$ of some $k'/Q_{j+1}$, i.e.
$|k'Q_j-kQ_{j+1}|\le t\bigl(Q_{j+1}/j+Q_j/(j+1)\bigr)$. Using
\eqref{eq:rec}, the number of surviving children per parent is bounded
\emph{uniformly} by $B=\lfloor16t\rfloor+1$ (Lemma~\ref{lem:branch},
machine-checked). Hence the number $N_n$ of level-$n$ intervals obeys
$N_n\le (Q_J+1)\,B^{\,n-J}$, while their diameter is $\le 2t/(nQ_n)$. By
Lemma~\ref{lem:factorial} (machine-checked), $Q_n\ge n!$, so for any $s>0$ the
$s$-content is bounded by
\[
N_n\cdot\bigl(2t/(nQ_n)\bigr)^{s}\le (Q_J+1) B^{n-J}(2t)^{s}\,(n\cdot n!)^{-s}
\longrightarrow0,
\]
because $\log(n!)\sim n\log n$ dominates the linear exponent $n\log B$. Thus
$\dimH C_J=0$ for each $J$, and $E_c$ is a countable union of such sets, so
$\dimH E_c=0$. Finally $F=\bigcup_nE_n$ is a countable union, hence
$\dimH F=0$.
\end{proof}

\paragraph{Lean status of the assembly.} All steps are kernel-checked.
The surviving centers are encoded as a \texttt{Finset} recursion
(\texttt{survSet}), not as paths: $S_0=\{0,\dots,Q_J\}$ and
$S_{n+1}=\bigcup_{k\in S_n}\mathrm{branchSet}(k)$, with
$|S_n|\le(Q_J+1)B^n$ by \texttt{Finset.card\_biUnion\_le\_card\_mul} and
Lemma~\ref{lem:branch} (\texttt{card\_survSet\_le}). Membership of the center
sequence $\mathrm{round}(Q_{J+n}x)$ in $S_n$ is proved by induction via the
triangle inequality (\texttt{centerOf\_mem\_survSet}), giving the cover
(\texttt{tailSet\_cover}). The content bound
$|S_n|(2t/((J+n)Q_{J+n}))^s\to0$ uses $Q_j\ge j!$ and
\texttt{FloorSemiring.tendsto\_pow\_div\_factorial\_atTop} with $c=B^{1/s}$
(\texttt{tendsto\_contentBound}); no asymptotic notation is used. Feeding this
into \texttt{hausdorffMeasure\_le\_liminf\_sum} yields $\mu H^s=0$ for all
$s>0$ (\texttt{hausdorffMeasure\_tailSet\_zero}), hence $\dimH=0$
(\texttt{dimH\_tailSet\_zero}). Periodicity
$\|Q_j(x+m)\|=\|Q_jx\|$ (\texttt{apprSeq\_add\_intCast}) reduces general $x$ to
$[0,1]$; integer shifts are $1$-Lipschitz and countable unions preserve
dimension $0$ (\texttt{dimH\_iUnion}). The sub-threshold case $c<1/4$ is
\emph{not} used: a monotonicity reduction to $c'=\max(c,0)$ lets the covering
argument handle all $c$ uniformly, so the theorem needs no continued-fraction
axioms. \texttt{Audit.lean} exits $0$ with \texttt{sorryAx} count $0$.

\begin{remark}
The inequality $\|Q_jx\|\le j\|Q_jx\|$ gives the inclusion from the
weighted zero-set into the unweighted one. To see strictness explicitly,
the point constructed in \texttt{QuarterPointFinal.lean} has
$j\|Q_jx_0\|\to1/4$. Its unweighted error tends to zero, and
$x_0\notin L$ by the weighted classification. The combined witness is kernel-checked in \texttt{UnweightedCorrection.lean} using only the three Lean base axioms.
\end{remark}

\begin{theorem}[Weighted upper bound]\label{thm:upper}
For $0\le\alpha\le1$ and every finite $c\ge0$,
$\dimH\{x:\Lam_\alpha(x)\le c\}\le1-\alpha$.
\end{theorem}
\begin{proof}[Proof sketch]
Cover $\{\Lam_\alpha\le c\}$ by $\bigcup_J\bigcap_{j\ge J}\{x:\|Q_jx\|\le t/j^\alpha\}$,
$t>c$. Level-$j$ intervals have radius $t/(j^\alpha Q_j)$; a parent has at most
$K j^{1-\alpha}$ surviving children (using $Q_{j+1}\le(4j+3)Q_j$). Thus
$N_n\le C K^n (n!)^{1-\alpha}$ and diameters $\le2t/(n^\alpha n!)$, so the
$s$-content is $\le C' K^n (n!)^{1-\alpha-s}\to0$ for $s>1-\alpha$
(the analytic core $K^n(n!)^{-\beta}\to0$, $\beta=s-(1-\alpha)>0$, is
machine-checked, \texttt{WeightedPhase.content\_core\_tendsto}). Hausdorff's
covering criterion gives $\dimH\le s$; let $s\downarrow1-\alpha$.
\end{proof}

\begin{theorem}[Prescribed limits at subcritical weights]\label{conj:weight}
Let $0\le\alpha<1$. For $\alpha>0$ take any $c\ge0$; for $\alpha=0$
take $0\le c\le1/2$. Then
\[
 \dimH\{x\in[0,1]: j^\alpha\|Q_jx\|\longrightarrow c\}=1-\alpha .
\]
In particular $\dimH\{x:\Lam_\alpha(x)\le c\}=1-\alpha$ for every
finite $c\ge0$, including $c=0$.
\end{theorem}
\begin{proof}
The upper bound follows from Theorem~\ref{thm:upper}. For the lower
bound put $w_j=j^{-\alpha}/\log(j+2)$ and use intervals
of radius $r_j=w_j/Q_j$ centred at $(k+cj^{-\alpha})/Q_j$, $k\in\Z$.
At all large levels $cj^{-\alpha}\in[0,1/2]$. The distance to the
integers is $1$-Lipschitz, so every surviving point satisfies
$|j^\alpha\|Q_jx\|-c|\le 1/\log(j+2)$, including at the endpoints.
Inside any parent interval, shrinking by $r_{j+1}$ at both ends leaves
at least $2(r_j-r_{j+1})Q_{j+1}-1$ eligible child centres, irrespective
of the shift. Since $Q_{j+1}/Q_j=4j+O(1)$, this is at least
$b_j=\lfloor j^{1-\alpha}/\log(j+2)\rfloor\ge2$ eventually.
Start with one parent inside $(0,1)$, keep exactly $b_j$ whole children
and split its mass evenly. The level-$n$ intervals are disjoint and
each has mass $M_n^{-1}$, with
$\log M_n=\sum_{j=N}^{n-1}\log b_j=(1-\alpha)n\log n+o(n\log n)$.
An interval of length $Q_{n+1}^{-1}\le h<Q_n^{-1}$ meets at most three
level-$n$ intervals by their grid spacing. Since
$\log Q_{n+1}=n\log n+O(n)$, its mass is at most $3/M_n\le h^s$
eventually for every $0<s<1-\alpha$. Summation over any sufficiently
fine cover proves the lower bound. For the sublevel statement take
the subset with limit zero and use the upper bound and integer translations.
\end{proof}

For $\alpha=0$ the dimension assertion follows from Corollary~8 of
\cite{BLR} with constant target $c$, by the distance inequality used
above. For $0<\alpha<1$, condition (1.6) of their Theorem~7 fails:
$\log Q_{j+1}/\log Q_j\to1$. The lower bound is also a corollary of
Feng, Wen and Wu's homogeneous Moran bounds (Theorem~2.1,
\emph{Sci. China Ser. A} 40 (1997), 475--482): the retained level
has $\log M_n\sim(1-\alpha)n\log n$ and both bounding logarithmic
length scales are asymptotic to $n\log n$. The independent covering
upper bound applies to the full prescribed-limit set.
\subsection{A Hausdorff gauge at the critical weight}
For an integer $k\ge1$, set $b=2k+1$. Define an increasing continuous
gauge $h_k$ by $h_k(0)=0$, $h_k(1/Q_n)=b^{-n}$ ($n\ge1$), and
linear interpolation on each intervening interval. Only its values
near zero matter. Write $\mathcal H^{h_k}$ for Hausdorff measure with
covering cost $h_k(\operatorname{diam} U)$.

\begin{theorem}[Size of a critical spectral level]\label{thm:gauge}
On the circle $\R/\Z$, the exact level $\{x:\Lam(x)=k/4\}$ has
infinite but $\sigma$-finite $\mathcal H^{h_k}$ measure.
The set $\{x:\Lam(x)<k/4\}$ has zero $\mathcal H^{h_k}$ measure.
\end{theorem}
\begin{proof}
Write $A_j=4j+2$, let $M_j$ be a nearest integer to $Q_jx$, put
$u_j=Q_jx-M_j$, and define $d_j=M_{j+1}-A_jM_j-M_{j-1}$.
Then $d_j=A_ju_j+u_{j-1}-u_{j+1}$.
If $\Lam(x)\le k/4$, eventually $|u_j|\le (k+1)/j$, and
$\limsup |d_j|\le k$. Thus the integer digits eventually satisfy
$|d_j|\le k$. For $\Lam(x)<k/4$ this improves to $|d_j|\le k-1$.
Partition into countably many classes by a common starting index $N$
and the two starting integers. By level $n$ there are at most
$b^{n-N}$ possible $M_n$ per class. Each gives an interval of diameter
$2(k+1)/(nQ_n)\le1/Q_n$ eventually. This gives a finite covering
cost per class; for the strict sublevel the cost tends to zero since
the count uses $2k-1<b$ instead. The classes may be taken Borel
by requiring the stated inequalities at all indices past $N$.

For positivity fix a large $N$, set $M_{N-1}=M_N=0$, and prescribe
arbitrary digits $d_j\in\{-k,\ldots,k\}$ for $j\ge N$ using the
recurrence above. With $W_j=M_{j+1}Q_j-M_jQ_{j+1}$ one has
$W_j=d_j Q_j-W_{j-1}$ and $W_{N-1}=0$. Since the denominators at
least double, $|W_j|\le2kQ_j$. Telescoping therefore gives a limit
$x=\lim M_j/Q_j$ with $|x-M_j/Q_j|\le4k/Q_{j+1}$ uniformly.
Choose $N$ so that $|u_j|<1/8$ for all $j\ge N$ and
$8k/Q_{N+1}<1/2$. Since $M_N=0$, the entire constructed set
lies in an interval of length at most $8k/Q_{N+1}$. On that
interval the projection to the circle preserves distances.
Thus $M_j$ is indeed nearest, and
$|4ju_j-d_j|\le2|u_j|+|u_{j-1}|+|u_{j+1}|\to0$.
The compact image of the digit space lies in the sublevel.

Give the digits equal independent weights. An interval of length
$1/Q_{n+1}\le h<1/Q_n$ determines at most one $M_j$ at each
$N\le j<n$: for two points in it their integers differ by less
than $1/4+Q_j/Q_n<1$. At level $n$ there are at most three choices.
These integers recover all digits through $n-1$, hence the interval
has mass at most $3b^{-(n-N)}\le3b^{N+1}h_k(h)$.
The mass distribution bound gives positive Hausdorff measure.
The local isometry transfers this bound to the circle.
Removing the null strict sublevel preserves positivity.

Choose a fixed rounding rule at ties. The functions
$x\mapsto j\|Q_jx\|$ are continuous on the circle, so their limsup
and its level sets are Borel. The digit image in the non-strict
sublevel retains positive measure upon removing the strict sublevel,
whose zero measure was proved independently by the smaller alphabet.
The countable covering classes above give sigma-finiteness also on
the exact level by intersection. For comparison, Cabrelli, Mendivil,
Molter and Shonkwiler~\cite[Thm.~3]{CMMS}
treat the gauge of a gap-generated Cantor set. We use the recurrence's
own scale and prove the assertion for the entire level here.
Finally, the exact level is a Borel set of Lebesgue measure zero by
Theorem~\ref{thm:finite}. For $y\in L$, the reverse triangle inequality gives
$|j\|Q_j(x+y)\|-j\|Q_jx\||\le j\|Q_jy\|\to0$.
Hence the exact level is invariant under translation by $L$ modulo one. This subgroup is dense as
$(e+1)/2$ is irrational. If its level measure were finite and positive,
restricting $\mathcal H^{h_k}$ to the level and normalizing would give
a probability invariant under this dense subgroup. Integrating continuous
functions and taking limits gives invariance under the entire circle.
Such a probability is normalized Lebesgue measure (equal arcs of length
$1/m$ have mass $1/m$, and rational arcs determine the measure).
It cannot be supported on the null level. This proves infinitude.
\end{proof}

\section{Numerical evidence}\label{sec:num}

All computations use mpmath with dynamic precision $\mathrm{digits}(Q_j)+300$
(to survive the catastrophic cancellation in $Q_jx-\mathrm{round}(Q_jx)$).
These are corroborations, not proofs.

\paragraph{Spectral gap.} For $x=B\pi$ ($B=1..10$), $\sqrt2$, the golden ratio,
and random reals, $\max_{j\le60}j\|Q_jx\|$ lies in $22$--$29$ and grows
$\sim j/4$, in this finite sample; the data do not establish an infinite limsup. For $x\in L$ (e.g.\ $(\ee+1)/2$,
$1+(\ee+1)/2$), $j\|Q_jx\|<10^{-12}$ already for $j\ge10$ and $\sim10^{-117}$
at $j=60$, confirming $\Lam=0$.

\paragraph{Dimension (box-counting).} Exact rational covering trees give the
following box-counting estimates for $E_c\cap[0,1]$:
\begin{center}
\begin{tabular}{c|cccc}
$c$ & $N=4$ & $N=12$ & path $N=400$ & bound $n=1500$\\\hline
$1/4$ & 0.371 & 0.323 & 0.171 & 0.180\\
$1/2$ & 0.630 & --- & 0.251 & 0.270\\
$1$   & 0.863 & --- & 0.343 & 0.360\\
$2$   & 1.000 & --- & 0.442 & 0.446
\end{tabular}
\end{center}
The finite counts suggest the branching scale $\lfloor8c\rfloor+1$ ($3,5,9,17$), and
the rigorous upper bound decays as $\log B/\log n\to0$. \textbf{Caveat:} the
decay is only $O(1/\log n)$; at accessible $N$ the estimates are far from $0$
(e.g.\ $s=0.01$ needs $n\approx10^{60}$). The numerics are consistent with
$\dimH=0$ but cannot resolve it directly; the proof is the covering argument of
Theorem~\ref{thm:finite}. Below the gap ($c=0.24$) the covering upper bound
still grows exponentially while the true branching tends to $0$: the covering
method is not sharp below the gap, which is why the exact classification
(Theorem~\ref{thm:class}) is needed there.

\section{Dimensions of continued-fraction window constraints}
Put $H_j=\lfloor Q_j\sqrt j\rfloor$. For irrational $x\in(0,1)$ let
$k(j)$ be its last convergent index with denominator at most $H_j$
(the integer part has index zero) and $r_j=k(j+1)-k(j)$.
\begin{theorem}\label{thm:windows}
For each prescribed sequence $\epsilon_j\in\{0,1\}$, the set with
$k(j)\equiv\epsilon_j\pmod2$ eventually has Hausdorff dimension one.
The set with eventually bounded $r_j$ has dimension $1/2$. Its subset
with $r_j=3$ and $k(j)+j$ odd eventually also has dimension $1/2$.
\end{theorem}
\begin{proof}
The recurrence gives $H_{j+1}/H_j\asymp j$ and
$\log H_j=j\log j+O(j)$. Use cylinders of finite regular continued
fractions. A cylinder with last two denominators $u,v$ has length
$1/(u(u+v))$. At stage $j$ stop at the first denominator crossing $H_j$,
with last digit four and $4H_j/3\le u\le 5H_j/2$.
To initialize at any sufficiently large stage choose the rational $1/d$
with $H_j/3\le d\le H_j/2$, use either terminal representation to
prescribe the index parity, and append four.

Write $T=H_{j+1}$, $R=T/u$. For the full-dimensional construction
choose coprime positive integers $m,n$ with
$0.34R\le m\le0.36R$ and $0.20R\le n\le0.22R$.
There are $\gg R^2$ choices uniformly for large $R$: by
M\"obius inversion, the count in this rectangle is its area divided
by $\zeta(2)$ plus $O(R\log R)$ (sum the two integer endpoint errors
over divisors up to $R$). Append the digits of $m/n$ as a finite tail.
Its terminal denominator is $t=mu+nv$, independently of the terminal
representation. Since $v\le u/4$, we have $T/3<t<T/2$.
Choose the representation with the desired terminal index parity and
append four. The new last denominator is between $4T/3$ and $5T/2$,
so the terminal index just chosen is exactly $k(j+1)$.
Distinct choices give distinct cylinders with disjoint interiors:
each word stops at its first crossing of $T$, and removing the final
four recovers the tail and its value.

For the lower bound with $r_j=3$, instead append $1,m,4$ with
$T/3\le m(u+v)+u\le T/2$. There are $\gg R$ choices and the first
crossing again occurs at the four. Initialize with $k(j)+j$ odd.
Three index steps preserve that condition at the next stage.

In either construction give siblings equal conditional masses.
With $b=2$ or $1$, respectively, the mass of a level-$j$ cylinder
is at most $\exp(-bj\log j+O(j))$. Such cylinders have disjoint
interiors and lengths comparable to $H_j^{-2}$ uniformly. An interval
of length $h\in[H_{j+1}^{-2},H_j^{-2}]$ meets at most a bounded number
of them. Since $\log H_{j+1}=j\log j+O(j)$, its mass is $O(h^s)$
for each $s<b/2$. The mass distribution principle proves the lower bounds.

Finally bounded increments imply $k(j)\le Cj$ eventually for some
integer $C$. Taking $j=\lfloor n/(2C)\rfloor$ gives $d_n>H_j$ and
$\log d_n/n\to\infty$. For $1/2<t<s$, summing the $t$-powers of
all level-$n$ cylinder lengths costs at most $\zeta(2t)^n$, since
$d_n\ge\prod_{i=1}^n a_i$. Restricting to $d_n\ge\exp(Mn)$ makes
the $s$-cost at most $(\zeta(2t)\exp(-2M(s-t)))^n$. Choose $M$
so the base is below one and take the countable union over starting
indices. This gives the upper bound $1/2$.
\end{proof}
For related spectra of continued-fraction growth see Fan, Liao, Wang
and Wu~\cite{FLLW}. The lower constructions here prescribe the
Hermite windows themselves. Removing the countable set $\Q+\Q e$
leaves both dimensions unchanged. These statements concern sets of
variable $x$ and leave the window pattern of the single point $\pi$
undetermined (translation by its integer part preserves denominators).

\section{Short continued-fraction windows occur infinitely often}
For $x\in(0,1)\setminus\mathbb Q$, retain $H_j,k(j),r_j$ from the window
constraints above.
\begin{theorem}[Short windows at almost every point]\label{thm:short-ae}
There is a constant $c>0$ such that for Lebesgue-almost every $x$,
\[
 \#\{1\le j\le J:r_j(x)\le1\}\ge c\log J
\]
for all sufficiently large $J$ (with a starting index depending on $x$).
\end{theorem}
\begin{proof}
Let $\tau_j=k(j)+1$ be the first index whose denominator exceeds $H_j$.
Stop each continued-fraction cylinder at this index and denote the
resulting sigma-field by $\mathcal F_j$. The stopping words are
prefix-free and cover the irrational points. Since $H_j$ increases,
$\mathcal F_j\subseteq\mathcal F_{j+1}$ and the event
$E_j=\{\tau_{j+1}-\tau_j\le1\}$ belongs to $\mathcal F_{j+1}$.

On a stopped cylinder write $q$ and $q'$ for the latest and previous
denominators, with $q>H_j$ and $0\le q'\le q$. If $q>H_{j+1}$, then
$\tau_{j+1}=\tau_j$. Otherwise put $A=\lfloor H_{j+1}/q\rfloor+1$.
The next digit $a\ge A$ ensures $a q+q'>H_{j+1}$ and gives $E_j$.
The conditional Lebesgue distribution of the remaining tail
$t=[0;a,\ldots]$ on the cylinder has density proportional to
$(q+q't)^{-2}$ for $0<t<1$. Thus
\[
 \mathbb P(a\ge A\mid\mathcal F_j)
 =\frac{q+q'}{Aq+q'}\ge\frac1A\ge \frac{H_j}{2H_{j+1}}.
\]
This bound also holds in the case $q>H_{j+1}$. The recurrence and the
definition of $H_j$ imply $H_{j+1}/H_j=O(j)$, so
$\mathbb P(E_j\mid\mathcal F_j)\ge c_0/j$ for all large $j$.

Write $Z_J=\sum_{j\le J}\mathbf 1_{E_j}$. The preceding indicators are
$\mathcal F_j$-measurable. Conditional expectation, iterated one step
at a time, gives
\[
 \mathbb E e^{-Z_J}\le C J^{-b},\qquad
 b=(1-e^{-1})c_0>0.
\]
In fact the factor at step $j$ is at most
$1-(1-e^{-1})c_0/j$. By Markov's inequality,
$\mathbb P(Z_J<(b/2)\log J)\le C J^{-b/2}$.
This is summable on $J=2^m$. The first Borel--Cantelli lemma and
monotonicity of $Z_J$ between successive powers of two prove the claim.
\end{proof}
The argument applies to any increasing integer thresholds tending to
infinity with successive ratios $O(j)$. The estimate concerns
Lebesgue measure on the variable $x$. It supplies a frequency bound
for the short-window criterion of the main paper.

\section{Eventual periodicity of digit residues}
For irrational $x=[0;a_1,a_2,\ldots]$ in $(0,1)$ let $E$ consist of
the points for which $(a_i\bmod4)$ is eventually periodic. Let $F$
consist of points with every digit in $\mathcal A=\{1,5,9,\ldots\}$.
Continued fractions with a fixed arithmetic-progression alphabet are
studied by Chousionis, Leykekhman and Urba\'nski~\cite{CLU}. Here
we compare different sequences of allowed residues, including sequences
that change from one position to the next.
\begin{theorem}\label{thm:residues}
$\dimH E=\dimH F$, with $13/20<\dimH F<5/7$.
\end{theorem}
\begin{proof}
Fix any prescribed sequence of residues. Add a fixed $c_i\in\{0,1,2,3\}$
to the $i$th digit of each point of $F$ to map it onto the desired set.
This map is Lipschitz uniformly in the sequence. Indeed, at the first
different digit $a<b$ the two source tails, written with integer part
zero, are at least $1/(a+1)-1/b$ apart. The new tails are at most
$1/a-1/(b+1)$ apart, since a common shift decreases this upper bound.
Their distance ratio is at most
\[
 \frac{b-a+1}{b-a-1}\frac{(a+1)b}{a(b+1)}\le10/3
\]
as $b-a\ge4$. A common prefix with last denominator $q$ has derivative
in absolute value between $1/(4q^2)$ and $1/q^2$. Increasing digits
can only increase the continuant $q$, so restoring the two prefixes
gives Lipschitz constant $40/3$. There are countably many eventually
periodic residue sequences, and $F\subset E$. This proves equality.

For the bounds write $f_a(x)=1/(a+x)$. In the bi-Lipschitz coordinate
$h_u(x)=1/(u+x)$, $u>0$, the absolute derivative of the conjugated
branch at $h_u(x)$ is
\[
 D_a(x)=\left(\frac{1+x/u}{a+x+1/u}\right)^2 .
\]
For the upper bound choose $u=5/2$, $s=5/7$ and $p=2s$.
Put $S(x)=\sum_{a\in\mathcal A}D_a(x)^s$.
After multiplication by a positive function of $x$, its derivative
has the sign of
\[
 1+1/u-u+\sum_{a\ge5,\ a\in\mathcal A}(a+1/u-u)
 \left(\frac{1+x+1/u}{a+x+1/u}\right)^{p+1}.
\]
Termwise differentiation is justified by uniform summability. This
expression increases with $x$, so the maximum of $S$ is at an endpoint.
Write $a=4k+1$. After the terms $k=0,1,2$, the integral from $2$
to infinity bounds the remaining sum by
$(7/12)(1+x/u)^p(9+x+1/u)^{-3/7}$.
At zero the three terms and this remainder are respectively less than
$(619,90,41,224)/1000$; at one they are less than
$(464,115,57,346)/1000$, by taking seventh powers. Thus $S\le0.982$.
By uniform convergence this bound stays strictly below one for some
$s<5/7$. Iterating the chain rule bounds the sum of powered derivatives
of all length-$n$ words at any one point by $(\sup S)^n$. Their
derivative distortion is uniformly bounded: ordinary prefix distortion
is at most four and $h_u$ and its inverse have derivatives bounded
above and away from zero. The cylinder $s$-cost therefore tends to zero.

For the lower bound use $u=2$. The infimum of $\sqrt{D_a(x)}$ on
$[0,1]$ is $3/5$ for $a=1$ and $1/(4k+3/2)$ for $a=4k+1$, $k\ge1$.
At $s=13/20$, using a decreasing-series integral from one to infinity,
\[
 \sum_{a\in\mathcal A}(\inf D_a)^s
 \ge(3/5)^{13/10}+(5/6)(11/2)^{-3/10}>1.
\]
The two terms exceed $51/100$ and $49/100$, respectively, by taking
tenth powers. A finite subalphabet therefore has sum $Z>1$ for some
$s>13/20$. Give its digits independent weights $(\inf D_a)^s/Z$.
By the chain rule a cylinder has mass at most its normalized length
to the power $s$ in the conjugated interval. Stop when its length
first drops below $h$. Uniform distortion and the finite alphabet
give lengths between $ch$ and $h$ for a fixed $c>0$. The stopped
intervals have disjoint interiors, so an interval of length $h$ meets
at most $2+1/c$ of them, up to endpoints of zero mass. The mass
distribution principle and bi-Lipschitz invariance prove the lower bound.
\end{proof}
The published bound of Chousionis, Leykekhman and Urba\'nski~\cite[p.~23, eq.~(4.8), arXiv version]{CLU}, applied to the alphabet $E_{1,4}$, also gives the upper bound $\dimH E<0.744$ by the equality above. The two bounds in the theorem follow from the branch estimates above.
In particular, removing $\mathbb Q-e$ from $E$ leaves its dimension
unchanged. For each remaining $x$, the sum $e+x$ is irrational even
though its digit residues pass the periodicity test applied to $x$.
The theorem gives no membership information for the fractional part of
$\pi$.
\section{Scope and verification}

\begin{itemize}[leftmargin=1.5em]
\item Theorem~\ref{thm:spectral}, Theorem~\ref{thm:sub}, and the countability /
dimension-$0$ statements for $c<1/4$ are represented in
\texttt{EpiFormal/SpectralGap.lean}. The declaration audit of those
results and their $\delta$ estimates reports only the three Lean axioms.
\item Theorem~\ref{thm:finite} (all $c\in\R$) is \emph{fully machine-checked}
(\texttt{DimHCore.lean} + \texttt{DimHFinite.lean}, base axioms only:
\texttt{propext}, \texttt{Classical.choice}, \texttt{Quot.sound}; no
$\delta$ axioms, no \texttt{sorryAx}). The assembly uses a surviving-center
\texttt{Finset} recursion rather than path encoding; the growth estimate is
explicit (\texttt{tendsto\_pow\_div\_factorial\_atTop}), not asymptotic.
\item The prescribed-limit and Hausdorff-gauge theorems, the window dimension theorem (Theorem~\ref{thm:windows}), and the residue comparison (Theorem~\ref{thm:residues}) have paper proofs.
\item This note does \emph{not} prove the irrationality of $e+\pi$ (still open).
It is a metric-theory spin-off; the connection is that the same classification
machinery arose from the $e+\pi$ escape program.
\item Section~\ref{sec:prior} describes the relation to earlier spectral-gap
and shrinking-target results for this choice of subsequence and weight.
\end{itemize}

\begin{thebibliography}{99}
\bibitem{CMMS} C.~Cabrelli, F.~Mendivil, U.~M.~Molter and R.~Shonkwiler,
\emph{On the Hausdorff $h$-measure of Cantor sets}, Pacific J. Math. 217
(2004), 45--59.
\bibitem{CLU} V.~Chousionis, D.~Leykekhman and M.~Urba\'nski, \emph{On the dimension spectrum of infinite subsystems of continued fractions}, arXiv:1805.11904.
\bibitem{FLLW} A.-H. Fan, L.-M. Liao, B.-W. Wang and J. Wu, \emph{On Khintchine exponents and Lyapunov exponents of continued fractions}, arXiv:0802.3433.
\bibitem{Phase1} Phase-1 archive, \emph{On the irrationality of $e+\pi$}
(formalization project), \texttt{EpiFormal/Classification.lean},
Theorem~2 and Corollary~3; \texttt{EpiFormal/SpectralGap.lean}, 2026-11-15.
\bibitem{KL} C.~Kraaikamp and P.~Liardet, \emph{Good approximations and
continued fractions}, Proc.~Amer.~Math.~Soc.~112 (1991), 303--309.
\bibitem{BDS} A.~B\'ir\'o, J.-M.~Deshouillers, V.~T.~S\'os, \emph{Good
approximation and characterization of subgroups of $\R/\Z$}, Studia Sci.~Math.~Hungar.~38 (2001), 97--113.
\bibitem{BLR} Y.~Bugeaud, L.~Liao, M.~Rams, \emph{Metrical results on the distribution of fractional parts of powers of real numbers}, Proc. Edinburgh Math. Soc. 62 (2019), 505--521, arXiv:1710.03596. (Cor.~8,
the hypothesis is $Q_{j+1}/Q_j\to\infty$.)
\bibitem{Moreira} C.~G.~Moreira, \emph{Geometric properties of the Markov and
Lagrange spectra}, arXiv:1612.05782.
\bibitem{LimaMoreira} D.~Lima and C.~G.~Moreira, \emph{Phase transitions for
the Lagrange and Markov spectra}, arXiv:1801.04636.
\bibitem{BV} V.~Beresnevich and S.~Velani, \emph{A mass transference principle
and the Duffin--Schaeffer conjecture}, Ann.~of Math.~164 (2006), 971--992.
\bibitem{Falconer} K.~Falconer, \emph{Fractal Geometry}, Wiley, 2003.
(countable sets have Hausdorff dimension $0$.)
\end{thebibliography}

\end{document}